\documentclass[11pt]{article}
% Shared preamble for per-Python-file documentation (input via \input{...}).
\usepackage[margin=1in]{geometry}
\usepackage[T1]{fontenc}
\usepackage[utf8]{inputenc}
\usepackage{lmodern}
\usepackage{hyperref}
\usepackage{xcolor}
\usepackage{listings}
\usepackage{listingsutf8}
\lstset{
  basicstyle=\ttfamily\small,
  columns=fullflexible,
  breaklines=true,
  upquote=true,
  inputencoding=utf8,
  extendedchars=true,
  frame=single,
  framerule=0.2pt,
  tabsize=2
}


\title{\texttt{run\_vr.py} (Q\&A)}
\author{}
\date{}

\begin{document}
\maketitle

\paragraph{Q: What is the purpose of \texttt{run\_vr.py}?}
\textbf{A:} It trains and evaluates a Wide \& Deep regression baseline (named ``VR'' here) for watch-time prediction. The key distinctive behavior is that the script rescales labels and predictions by a constant factor during training/evaluation.

\paragraph{Q: How does this script load data and create the model?}
\textbf{A:} It loads \lstinline|{dataset_path}/{dataset_name}/{dataset_name}_data.pkl| via \lstinline|KUAIRECDataLoader| and constructs \lstinline|WideAndDeep| with a configurable MLP size and embedding dimension.

\paragraph{Q: What is special about the label scaling in this file?}
\textbf{A:} During training, it trains with MSE against \lstinline|20 * label|, and during evaluation it divides predictions by 20 (\lstinline|y = model(features) / 20|). This is a numerical scaling choice (it does not change the optimum under MSE if done consistently).

\paragraph{Q: What metrics are printed at the end?}
\textbf{A:} It computes MAE/XAUC/KL using \lstinline|utils.py| and rescales MAE to seconds for the configured dataset.

\paragraph{Q: Code example: the scaling logic?}
\textbf{A:}
\begin{lstlisting}[language=Python]
# train
loss = criterion(y, 20 * label.float())

# test
y = model(features) / 20
\end{lstlisting}

\paragraph{Q: Full source code?}
\textbf{A:}
\lstinputlisting[language=Python]{/workspace/run_vr.py}

\end{document}
